
\def\PUDcodigo{EME125}

\def\PUDcodacad{04506.44}

\def\PUDdisciplina{Análise de Falhas}

\def\PUDsemestre{S9}

\def\PUDhoras{80}

%NOVOS ---------------------------------------------------
\def\PUDhorasTeoria{40}
\def\PUDhorasPratica{40}

\def\PUDinterECA{
\hyperref[PUD:EME112]{Soldagem (04506.29)}

\hyperref[PUD:EME119]{Sistemas Mecânicos (04506.35)}

\hyperref[PUD:EME128]{Inspeção Industrial (04506.50)}

\hyperref[PUD:EME102]{Tribologia e Lubrificação (04506.53)}
}

\def\PUDatuaisECA{---}

\def\PUDrecursos{Projetor multimídia; Quadro branco e pincel; Aulas em laboratório.}

%----------------------------------------------------------

\def\PUDcreditos{4}

\def\PUDprerequisito{ \hyperref[PUD:EME109]{EME109} }

\def\PUDpreacad{ \hyperref[PUD:EME109]{Resistência dos Materiais I (04506.26)} }

\def\PUDobjetivos{Conhecer os principais modos de falhas em componentes de equipamentos;  Reconhecer as características dos principais modos de falhas;  Estimar tempo de vida de componentes de equipamentos;  Projetar componentes de equipamentos com relação à fratura;  Avaliar risco de falha em componentes de equipamentos industriais.}

\def\PUDementa{Origem e Fundamentos da Mecânica da Fratura;  Bases da Mecânica Linear da Fratura;  Elementos
da Mecânica Não-Linear da Fratura;  Fratura por Carregamento Cíclico;  Introdução à Mecânica do Dano Contínuo.}

\def\PUDconteudo{  & Modos de falha: Deformação elástica, escoamento, ruptura dúctil, fratura frágil, fadiga, corrosão, desgaste, impacto, fretagem, fluência, flambagem, etc. 
\\ 
& Origem e fundamentos da mecânica da fratura: Termos e conceitos iniciais, notações e parâmetros da mecânica dos sólidos, comportamento mecânico dos materiais estruturais, resistência teórica dos materiais e concentração de tensões. 
\\ 
 & Introdução à mecânica do dano contínuo: Parâmetro escalar do dano, conceitos avançados, propagação subcrítica de trincas e tensões nas vizinhanças da ponta da trinca. 
\\ 
 & Bases da mecânica linear da fratura: Conceito de Griffith, distribuições de tensões na ponta da trinca, modo de cisalhamento antiplano, modo de tração e modo de cisalhamento plano, tendacidade à fratura e fratura quase frágil. 
\\ 
 & Elementos da mecânica não-linear da fratura: Zona plástica nas vizinhanças da ponta da trinca, modelo de trinca com zona plástica fina, critério deformacional da fratura, integral “J” e critérios energéticos. 
\\ 
 & Metodologia de análise de falha: Análise inicial da falha, levantamento de dados, análise da fratura, definição dos testes a serem realizados, interpretação dos resultados, relatório de análise de falhas. 
%\\ 
% & Técnicas de análise e diagnóstico: Inspeção visual, análise química, ensaios mecânicos, microscopia, etc. 
}

\def\PUDmetodo{Aulas expositórias que podem ser teóricas e/ou práticas, onde as práticas no laboratório serão marcadas no decorrer da disciplina.}

\def\PUDavalia{A avaliação será feita com aplicação de provas teóricas e/ou práticas, além da possibilidade de inclusão de trabalhos e seminários no decorrer da disciplina.}

\def\PUDbibbasica{ & \href{www.biblioteca.ifce.edu.br/index.asp?codigo_sophia=61962}{Pelliccione}, A. S.; Moraes, M. F.; Galvão, J. L. R.; Mello, L. A.; Silva, E. S.  Análise de Falhas em Equipamentos de Processo - Mecanismos de Danos e Casos Práticos.  2º ed.Rio de Janeiro, RJ: Interciência, 2014.  386p.  ISBN: 9788571933286. 
\\ 
 & \href{www.biblioteca.ifce.edu.br/index.asp?codigo_sophia=23902}{Collins}, J.  Projeto Mecânico de Elementos de Máquinas: uma perspectiva de prevenção da falha.  Rio de Janeiro, RJ: LTC, 2006.  740p.  ISBN: 8521614756. 
\\ 
& \href{www.biblioteca.ifce.edu.br/index.asp?codigo_sophia=63255}{Rocha}, J. A. L.  Termodinâmica da Fratura – Uma Nova Abordagem do Problema da Fratura nos Sólidos.  Salvador.  Editora EDUFBA.  2010.  195p.  ISBN: 9788523206956.
 %& PASTOUKHOV, V.  A.;  VOORWALD, H.  J.  C.;  Introdução à Mecânica da Integridade Estrutural.  1ª edição.  São Paulo.  Editora UNESP.  1995.  196p.  ISBN 9788571390805.  
 %\href{www.biblioteca.ifce.edu.br/index.asp?codigo_sophia=44596}{Garcia}, A.; Spim, J. A.; Santos, C. A.  Ensaios dos Materiais.  2ª ed.  Rio de Janeiro.  Editora LTC.  2012.  365p.  ISBN: 9788521620679. 
 }

\def\PUDbibcomplementar{ &  \href{biblioteca.ifce.edu.br/index.asp?codigo_sophia=62087}{Affonso}, Luiz Otávio Amaral. Equipamentos mecânicos: análise de falhas e solução de problemas. 3. ed. Rio de Janeiro: Qualitymark, 2014. 387 p., il. ISBN 9788541400367
\\
 & \href{biblioteca.ifce.edu.br/index.asp?codigo_sophia=44596}{Garcia}, Amauri; SPIM, Jaime Alvares; SANTOS, Carlos Alexandre dos. Ensaios dos materiais. 2. ed. Rio de Janeiro: LTC, 2012. 365 p. ISBN 9788521620679.
 \\
& \href{www.biblioteca.ifce.edu.br/index.asp?codigo_sophia=59389}{Serra}, Eduardo Torres.  Análise de falhas em materiais utilizados no setor elétrico.  E-book.  Interciênca.  554p.  ISBN: 9788571933712.
\\
 & \href{www.biblioteca.ifce.edu.br/index.asp?codigo_sophia=23805}{Souza}, Sérgio Augusto de.  Ensaios mecânicos de materiais metálicos: fundamentos teóricos e práticos.  5º ed.  São Paulo, SP: Blucher, 2011.  286p.  ISBN: 9788521200123.
\\
 &  \href{biblioteca.ifce.edu.br/index.asp?codigo_sophia=24118}{Colpaert}, Hubertus. Metalografia dos produtos siderúrgicos comuns. 4. ed. São Paulo: Edgard Blücher, 2008. 652 p. + Inclui CD-ROM. Inclui bibliografia e índice. ISBN 9788521204497. }

\def\PUDautor{Venceslau Xavier de Lima Filho}

\def\PUDdata{2013-06-17}

\def\PUDrevisao{2}

\def\PUDrevisor{Venceslau Xavier de Lima Filho}

\def\PUDdatarevisao{2019-05-15}

\PUDcorpo
